% Section 08: Discussion
% Mode: Interpretive — "These findings suggest...", "In our view...",
%   "This aligns with \textcite{author}'s concept of..."

\section{Discussion}\label{sec:discussion}

\subsection*{Introduction}

This section interprets the findings from \textit{Results and Analysis} (Section~\ref{sec:results}) and the
guidance artifact from \textit{Guidance Artifact} (Section~\ref{sec:artifact}) against the theoretical
framework established in the \textit{Theoretical Background} (Section~\ref{sec:theoretical-background})
and \textit{Conceptual Model} (Section~\ref{sec:conceptual-model}). Three lenses structure the interpretation:
distributed cognition, practice theory, and the jagged technological frontier.
The section then triangulates the practice catalog against external evidence,
draws implications for Enterprise Architecture (EA) practice and theory, and
closes with limitations and future research directions. The following
subsections detail this interpretation:

\localtableofcontents\clearpage

%% ================================================================
%% 8.2 Distributed Cognition in EA-GenAI Practice
%% ================================================================
\subsection{Distributed Cognition in EA-GenAI Practice}\label{subsec:disc-dcog}

The conceptual model positioned EA-GenAI integration as a distributed cognitive
system where the functional properties of the whole cannot be reduced to the
properties of individual agents \parencite{hutchins_cognition_1995}. The
Delphi findings confirm this
framing while revealing specific distribution patterns that the theoretical
framework anticipated but did not specify.

\subsubsection{Cognitive Distribution Patterns Across Practices}

The seven validated practice categories represent distinct configurations of
cognitive distribution between human and artificial agents. These
configurations align with \textcite{hollan_distributed_2000}'s three types of
distribution (across group members, between internal and external structures,
and through time) though the boundaries are more fluid than the typology
implies.

The Delphi data suggest that System/Codebase Analysis distributes information
retrieval to GenAI while retaining interpretive judgment with the practitioner.
The tight consensus on this practice
(Table~\ref{tab:practice-importance}) suggests that cognitive
distribution succeeds most clearly when the delegated task (scanning code
repositories, mapping dependencies) is verifiable against physical systems.
The trust spectrum for GenAI in RA development (Table~\ref{tab:trust-spectrum}) places information
compression at the highest trust level, consistent with this interpretation.
This convergence implies that verifiability, not task complexity, determines
where cognitive distribution succeeds most readily.

Research \& Exploration appears to distribute information gathering across a
broader knowledge space. The prevalence of Research Assistant mode among Round~1
practices (Section~\ref{subsubsec:practice-discovery}) suggests that
practitioners default to using GenAI to compress
large bodies of information into actionable summaries. This corresponds to
\textcite{hollan_distributed_2000}'s coordination between internal and external
structures: the practitioner's architectural mental model interacts with
GenAI's encoded knowledge to produce synthesis neither could achieve alone.

In a similar vein, First Draft \& Documentation appears to distribute
expression while retaining architectural thinking. The widespread adoption of
the Compression Not Generation framing (Section~\ref{subsubsec:cross-cutting})
captures this distinction precisely. As P3 stated: ``The thinking is still mine; the
AI accelerates expression'' (R1 Workbook). GenAI handles the cognitive labour
of transforming ideas into structured documents; the practitioner retains the
conceptual architecture that gives those documents meaning.

Stakeholder Communication and Governance/Compliance occupy a different position
in the distribution landscape. Both exhibited high variability
(Table~\ref{tab:practice-importance}), and both involve organizational context
that the panel near-unanimously identified as outside GenAI capability
(Section~\ref{subsubsec:cross-cutting}). These findings suggest that cognitive
distribution encounters a boundary at the organizational context interface; a
point we return to below.

\subsubsection{The Organizational Context Boundary}

The most consistent finding across the Delphi study was the organizational
context boundary. The panel near-unanimously identified team dynamics,
political constraints, and historical decisions as knowledge that cannot be
delegated to GenAI (Section~\ref{subsubsec:cross-cutting}). P3 observed that
``organisational context is nearly impossible to convey'' (R1 Workbook), while
P4 noted that GenAI ``is not valuable for the sociotechnical aspects that
determine whether an architecture will actually succeed'' (R1 Workbook).

In our view, this boundary maps directly to
\textcite{hutchins_cognition_1995}'s concept of functional system boundaries.
A distributed cognitive system operates within constraints set by its
participants' capabilities and the media available for coordination.
Organizational context (tacit, political, embedded in history) resists being
translated into forms GenAI can process. The context engineering practices described in
Section~\ref{subsubsec:context-engineering} narrow this boundary by making more
organizational knowledge machine-accessible, but they do not eliminate it.

This interpretation aligns with \textcite{saint-louis_investigation_2019}'s
three contexts of EA practice. The technological context maps readily to
cognitive distribution (System/Codebase Analysis). The socio-technological
context permits partial distribution with human oversight (First Draft \&
Documentation, Adversarial Review). The eco-technological context (external
standards, compliance, regulatory relationships) demands the most careful
verification (Governance/Compliance). An alternative explanation is that this boundary reflects current GenAI
limitations rather than a structural feature; future systems with richer
organizational memory may narrow it. However, the consistency across
all three EA contexts suggests the boundary reflects a structural
feature of distributed cognition in sociotechnical systems, not solely a
technological shortcoming.

\subsubsection{Verification as Cognitive Coordination}

The universal verification mandate (all eight participants described explicit
verification practices, with none accepting GenAI outputs without review
(Section~\ref{subsubsec:cross-cutting})) takes on new significance when viewed
through the distributed cognition lens. Verification is not merely quality
control. It is the mechanism through which the distributed cognitive system
maintains coherence.

\textcite{hutchins_cognition_1995} described how navigation teams use
redundant representations and cross-checking procedures to detect and correct
errors propagating through the cognitive system. The Delphi panel's
verification practices serve an analogous function. The top-tier heuristics (cross-validation, incremental trust building, and
context management; Table~\ref{tab:heuristic-importance}) are cognitive
coordination mechanisms that ensure distributed outputs align with the
practitioner's architectural intent.

The verification burden shift, rated the most significant practice change
(Table~\ref{tab:change-significance}), reflects a
fundamental restructuring of the cognitive system. Practitioners spend less
time on information gathering and expression (now distributed to GenAI) and
more time on coordination and validation. This shift is consistent with
\textcite{hollan_distributed_2000}'s observation that distributed cognitive
systems require coordination overhead proportional to their distribution
complexity.

%% ================================================================
%% 8.3 Practice Evolution Through a Practice Theory Lens
%% ================================================================
\subsection{Practice Evolution Through a Practice Theory Lens}\label{subsec:disc-practice-theory}

The conceptual model drew on \textcite{feldman_theorizing_2011}'s
conceptualization of practice as recurrent, recognizable patterns of action that
emerge through repeated performance. The Delphi findings provide empirical grounding for how
such patterns emerge under conditions of rapid technological change.

\subsubsection{From Creator to Curator: A Sociomaterial Accomplishment}

The creator-to-curator role transformation
(Section~\ref{subsubsec:practice-changes}) represents what
\textcite{orlikowski_sociomaterial_2007} would describe as a sociomaterial
accomplishment. The role shift did not arise from a deliberate redesign of
practice. It emerged through situated interaction between practitioners and
GenAI capabilities.

\textcite{orlikowski_using_2000}'s concept of technology-in-practice
illuminates this process. Practitioners did not adopt GenAI according to a
predefined workflow. Instead, they enacted new practices through ongoing
interaction, discovering that their most productive contribution shifted from
producing architectural content to evaluating, editing, and validating it.
Round~2 confirmed this shift: verification burden received the highest
significance rating, while the creator-to-curator transformation showed
notably higher variability (Table~\ref{tab:change-significance}).

A selection bias caveat applies: practitioners who gravitate toward curation
may be those who adopted GenAI earliest. However, the variability in the
creator-to-curator rating is itself informative.
\textcite{feldman_theorizing_2011} emphasize that practices are constituted
through performance. The uneven adoption of the curator role across the panel
suggests that this transformation is still being constituted, with different
organizational contexts producing different enactments.

\subsubsection{Emergent Practice Patterns}

The seven practice categories exhibit the characteristics
\textcite{feldman_theorizing_2011} identify as constitutive of practice:
recurrence, recognizability, and emergence through repeated performance. None
of the 19 Round~1 practices were formally designed. All arose through
practitioner experimentation with GenAI capabilities in the context of
architectural work. This emergence aligns with the conceptual model's
prediction that ``practices cannot be predicted from GenAI's technical
capabilities alone'' (Section~\ref{sec:conceptual-model}).

The Delphi convergence process itself demonstrated how individual practices
aggregate into collective patterns. Round~1 captured 19 individual practices;
Round~2 synthesized these into seven validated categories; Round~3 refined
heuristics and cross-role dynamics. This progression mirrors the
individual-collective dimension identified in the conceptual framework
(Section~\ref{subsec:frameworkdimensions}): personal experimentation aggregated
into shared organizational patterns through collective deliberation.

\textcite{kotusev_enterprise_2018}'s finding that successful EA practices
diverge substantially from mainstream literature prescriptions resonates
with this emergence pattern. The seven practice categories did not map neatly
onto existing EA frameworks or prescribed GenAI workflows. They emerged from
the complex interplay of practitioner creativity, organizational constraints,
and technological capabilities; precisely the sociomaterial dynamics that
practice theory predicts.

The cross-role dynamics data reinforce this interpretation. All six cross-role
patterns in Table~\ref{tab:cross-role} exhibited variability between 37\% and
79\% (Section~\ref{subsubsec:cross-role}), indicating that practices crossing
traditional role boundaries are still actively being constituted. The ``Ghost
Architect'' provocation surfaced a genuine tension: GenAI democratized artifact
production, yet accountability for architectural decisions remained with the
architect. This asymmetry (distributed production with
concentrated accountability) illustrates the uneven pace at which different
dimensions of practice emerge. \textcite{feldman_theorizing_2011} argue that
practices are constituted through performance; when new actors enter the
performance space, the practice itself must be reconstituted. The high
variability across cross-role patterns suggests that this reconstitution is
underway but far from settled, shaped more by organizational context than by
GenAI capability alone.

\subsubsection{Compression Not Generation as Practice Reframing}

The Compression Not Generation pattern represents a collective reframing of
GenAI's role in architectural practice, adopted by the majority of the panel. As P2 described: ``The value is not in generation
but in compression. It compresses reading time for surveys'' (R1 Workbook).
This framing is significant from a practice theory perspective because it
reveals how practitioners make sense of a new technology in terms compatible
with their existing professional identity.

\textcite{orlikowski_technological_1994}'s concept of technological frames
helps explain this dynamic. By framing GenAI as a compression tool rather than
a generative agent, practitioners maintained their self-understanding as the
source of architectural knowledge. The technology became an accelerator of
human thinking rather than a substitute for it. This framing connects directly
to the role identity preservation pattern: all eight participants viewed GenAI
as an efficiency tool, not a replacement for professional judgment
(Section~\ref{subsubsec:cross-cutting}).

%% ================================================================
%% 8.4 The Jagged Frontier in Architectural Work
%% ================================================================
\subsection{The Jagged Frontier in Architectural Work}\label{subsec:disc-jagged-frontier}

\textcite{dellacqua_navigating_2023} introduced the concept of a ``jagged
technological frontier'' to describe how AI capabilities vary unpredictably
across tasks. Their study of Boston Consulting Group consultants found that
task performance improved substantially when tasks fell inside the frontier
but degraded when they fell outside it. The Delphi consensus--variability
profiles map onto this concept in revealing ways.

\subsubsection{Mapping the Frontier Through Consensus Data}

System/Codebase Analysis sits clearly inside the frontier
(Table~\ref{tab:practice-importance}). The tight consensus and low variability
suggest that organizational context matters less: codebase analysis is
relatively standardized, verifiable, and bounded. This stable frontier boundary
suggests that directive guidance may be feasible.

Research \& Exploration and Adversarial Review occupy the frontier's middle
zone, with moderate consensus supporting contextual guidance, namely general
recommendations with situational qualification. The frontier here depends on
the specificity of the research question and the practitioner's domain
expertise.

Stakeholder Communication and Governance/Compliance sit on or outside the
frontier. Their high variability reflects that GenAI's value for these
practices depends heavily on organizational context; precisely the knowledge
that the panel identified as outside GenAI capability. The shifting frontier
boundary suggests that adaptive guidance may be more appropriate for these
practices.

Deep Workflow Integration presents a different pattern. Its low mean rating
and one abstention likely reflect adoption barriers rather than low frontier
placement. The participants operating in Cyborg mode
(Section~\ref{subsubsec:practice-discovery}) demonstrated that
practitioners who overcome infrastructure prerequisites achieve sustained
workflow benefits. The frontier for this practice may be wide but the path to reaching it is
steep. This suggests that for Deep Workflow Integration, the primary barrier
is organizational readiness rather than GenAI capability.

\subsubsection{Frontier Learning Through Failure}

The majority of participants described learning GenAI's limitations through
errors: hallucinated benchmarks, non-existent folder references, bounded
context suggestions that were ``technically logical but organisationally
nonsensical'' (P3, R1 Workbook). These incidents demonstrate what
\textcite{dellacqua_navigating_2023} describe as the cost of operating outside
the frontier: degraded performance that can propagate downstream if
undetected.

In our view, the frontier learning pattern reveals a distinctive feature of the
EA-GenAI cognitive system. Unlike the BCG consultants in
\textcite{dellacqua_navigating_2023}'s study, who were experimentally assigned
to tasks inside or outside the frontier, EA practitioners must navigate the
frontier continuously in their daily work. The verification heuristics refined
across three Delphi rounds (Table~\ref{tab:heuristic-importance}) represent
collectively developed navigation strategies for this frontier.

The tier structure of verification heuristics reinforces this interpretation.
Process-oriented heuristics (cross-validation, incremental trust building,
context management) rated higher than experience-dependent ones (domain
expertise threshold, learn through failure). This suggests that the panel
valued systematic frontier navigation strategies over individual intuition; a
finding consistent with \textcite{hutchins_cognition_1995}'s emphasis on
procedural coordination in distributed cognitive systems.

\subsubsection{Trust Calibration as Frontier Navigation}

Trust calibration (the bidirectional process of adjusting expectations about
GenAI capabilities) functions as the primary mechanism for navigating the
jagged frontier. The ``junior colleague'' mental model recurred across the
Round~3 panel: ``Think of AI as a junior architect and yourself as the senior
architect'' (R3 Brainstorm). This framing establishes a calibration anchor that
adapts as practitioners accumulate experience with GenAI's capability profile.

The panel confirmed that trust calibration is teachable but requires domain
expertise as a prerequisite (Section~\ref{subsubsec:teachability}). This
finding has implications for the frontier's temporal dimension. Novice
practitioners lack the domain knowledge to recognize when GenAI operates
outside the frontier, making them more vulnerable to the ``beautiful but wrong''
outputs that \textcite{dellacqua_navigating_2023} identified. The experience
indicators in the guidance artifact (Section~\ref{subsec:practice-catalog})
address this by providing graduated entry points for practitioners at different
experience levels.

%% ================================================================
%% 8.5 External Triangulation of Delphi Practices
%% ================================================================
\subsection{External Triangulation of Delphi Practices}\label{subsec:disc-triangulation}

The guidance artifact in Section~\ref{sec:artifact} is grounded in the Delphi
study. This section examines how external evidence corroborates, extends, or
qualifies the seven validated practice categories. The purpose is triangulation:
establishing whether the Delphi findings reflect broader practitioner experience
or remain idiosyncratic to this panel.

\textbf{System/Codebase Analysis.} Practitioner evidence confirms codebase
analysis as the highest-confidence GenAI application. Practitioners report using
GenAI for reverse-engineering software architecture, producing dependency graphs
and cross-repository documentation
\parencite{chatlatanagulchai_agent_readmes_2025}. Agent context files containing
architecture content appear in 67.7\% of analyzed repositories, indicating
widespread adoption of this practice in tooling-integrated workflows. This
convergence strengthens the directive guidance level assigned to this practice.

\textbf{Research \& Exploration.} Multi-model research pipelines are emerging,
where practitioners use different AI models for planning, structuring, and
execution stages \parencite{banh_affordance_genai_ea_2025}. The compression
framing (GenAI compresses existing information rather than creating new
knowledge) is well-established in practitioner discourse and aligns with the
Delphi finding that five of eight participants adopted this framing
(Section~\ref{subsubsec:cross-cutting}). External evidence thus corroborates
both the practice and its underlying cognitive model.

\textbf{Adversarial Review / Stress Testing.} The adversarial use of GenAI is
documented across practitioner discourse, though primarily as informal validation
rather than structured architectural critique
\parencite{lee_devils_advocate_2025, stanimirovic_ai_sparring_2026}. Epistemic
stress testing (tracking confidence levels in AI-mediated architectural
decisions) connects to the distributed cognition lens by making cognitive
assumptions explicit \parencite{gilda_epistemic_status_2026}. The stress testing
techniques identified in Table~\ref{tab:stress-testing} draw on both Delphi data
and this external evidence.

\textbf{Governance/Compliance.} The gap between governance aspiration and
implementation is well-documented.
\textcite{ettinger_ea_dynamic_capability_2025} identifies significant
limitations in existing EA frameworks for addressing GenAI requirements, noting
tension between fostering innovation and maintaining compliance.
Architecture-as-Code approaches (encoding architectural constraints as
machine-readable rules in CI/CD pipelines) represent an emerging response to this
gap \parencite{bucaioni_architecture_as_code_2025}. The high variability in the Delphi rating
(Table~\ref{tab:practice-importance}) is consistent with this literature:
governance value depends heavily on organizational maturity.

\textbf{First Draft \& Documentation.} This practice has the strongest external
corroboration. Practitioners report generating architecture decision records from
meeting transcripts and requirements, with multi-layer verification workflows
combining metaprompts, fresh LLM contexts, and LLM-as-Judge evaluation before
human review \parencite{banh_affordance_genai_ea_2025}. The verification burden
shift (the highest-rated change theme,
Table~\ref{tab:change-significance}) is most directly experienced in
documentation workflows. A related emerging practice, Spec-Driven Development
(SDD), extends the documentation-first principle by using formal specifications
as the primary human-authored artifact from which implementations derive
\parencite{piskala_spec_driven_dev_2026}. SDD is not yet a validated Delphi
practice, but its specification-anchored approach mirrors how RAs function:
normative specifications constraining downstream implementation. Early evidence
suggests three maturity levels (spec-first, spec-anchored, and spec-as-source),
though concerns about specification drift and tool inflexibility remain
\parencite[sec.~IX]{piskala_spec_driven_dev_2026}.

\textbf{Stakeholder Communication.} External evidence for GenAI-assisted
stakeholder communication in architecture contexts is sparse compared to other
practices. The strongest evidence concerns meeting transcript synthesis into
formal records (collective decision capture received the highest cross-role
pattern rating, Table~\ref{tab:cross-role}) rather than proactive
stakeholder communication. This evidence gap is consistent with the practice's
high variability and adaptive guidance level.

\textbf{Deep Workflow Integration.} The shift from chat-based interaction to
IDE-integrated workflows represents a maturation trajectory documented across
practitioner discourse \parencite{rajasekaran_context_engineering_2025,
chatlatanagulchai_agent_readmes_2025}. Agent context files containing
architecture content appear in 67.7\% of studied repositories, and enterprise
organizations are converting internal APIs into MCP servers at scale. Context
engineering (designing the information environment for AI agents) is emerging as
a core architectural skill (Section~\ref{subsubsec:context-engineering}).

Taken together, external evidence corroborates the Delphi findings across all
seven practice categories, though with varying strength. System/Codebase
Analysis and First Draft \& Documentation have the strongest external
validation, consistent with their higher consensus ratings. Stakeholder
Communication and Deep Workflow Integration have sparser external evidence,
consistent with their adaptive guidance levels. This pattern suggests that
the Delphi panel's consensus--variability profiles reflect genuine differences
in practice maturity rather than panel idiosyncrasy.

%% ================================================================
%% 8.6 Implications for EA Practice
%% ================================================================
\subsection{Implications for EA Practice}\label{subsec:disc-implications-practice}

The theoretical interpretations above yield four practical implications for
how EA practice may evolve as GenAI integration matures.

\subsubsection{Context Engineering as an Emerging Competency}\label{subsubsec:disc-context-eng}

The Delphi findings position context engineering (systematically designing
the information provided to GenAI systems) as an emerging EA competency.
The strong panel endorsement of the context management heuristic
(Table~\ref{tab:heuristic-importance}), combined with the finding that
agent context files containing architecture content appear in 67.7\% of
studied repositories \parencite{chatlatanagulchai_agent_readmes_2025},
indicates that context engineering is not merely aspirational but reflects
emerging practice maturity. These files function as both human-readable
documentation and machine-executable constraints.

Context engineering extends the traditional EA practitioner
competency set identified by \textcite{land_enterprise_2009} (which
emphasized technical expertise, communication skills, and political
acumen) with a new dimension. That dimension is the ability to represent
architectural knowledge in forms that GenAI systems can process
effectively. This is not
merely a technical skill. It requires the architectural judgment to determine
which constraints matter, the communication skill to encode them precisely,
and the organizational awareness to know what context the AI lacks.

\subsubsection{The Gate-to-Guardrail Transition}

The governance adaptation themes in Section~\ref{subsec:governance-adaptation}
suggest a shift from periodic gate reviews to continuous guardrail enforcement.
GenAI's acceleration of artifact production (practitioners producing drafts in
minutes rather than days) challenges the temporal assumptions of traditional
architectural review boards. The Delphi panel's low rating for inside-out governance
(Table~\ref{tab:cross-role}) suggests that most organizations have not yet
confronted this shift.

\textcite{ettinger_ea_dynamic_capability_2025} identifies significant
limitations in existing EA frameworks for addressing GenAI requirements. In our
view, the gate-to-guardrail transition represents the governance equivalent of
cognitive distribution: encoding architectural constraints as machine-readable
rules distributes governance enforcement across the development pipeline rather
than concentrating it in periodic human reviews. This transition directly
addresses the accountability gap exposed by the Ghost Architect dynamic
(Section~\ref{subsubsec:cross-role}): when artifact production is distributed
but accountability remains concentrated, continuous guardrail enforcement shifts
evaluation from \textit{who produced} the artifact to \textit{whether it meets
encoded standards}. In this way, the gate-to-guardrail model resolves the
asymmetry between distributed production and concentrated accountability by
making compliance verifiable regardless of the artifact's origin.

\subsubsection{The Competency Decay Paradox}

The skills being automated (cold-start writing, exhaustive manual search,
cover-to-cover documentation reading) are the same skills through which
practitioners built their expertise
\parencite{kabashkin_cognitive_atrophy_2025}. The panel rated ``skills used less'' as the second-most significant change
(Table~\ref{tab:change-significance}), confirming that practitioners recognize
this dynamic. Yet ``new skills emerging'' received a notably lower rating,
suggesting practitioners perceive the shift as extending existing expertise
rather than requiring fundamentally new competencies.

This tension reveals a paradox that practice theory helps illuminate. If
practices are constituted through performance
\parencite{feldman_theorizing_2011}, then automating the performance
undermines the mechanism through which expertise develops. Future practitioners
may lack the foundation required for effective trust calibration; precisely
the skill the panel identified as essential for all seven practices.

The anti-atrophy strategies proposed in
Section~\ref{subsubsec:anti-atrophy} constitute a practice-theoretic response
to this paradox. If expertise develops through performance
\parencite{feldman_theorizing_2011}, then preserving expertise requires
deliberately reintroducing the performance that automation displaces. Each
strategy targets a different facet of this mechanism. AI-free design sessions
preserve foundational skill performance by requiring practitioners to produce
architectural work without GenAI assistance. The attempt-independently-first
protocol maintains mental model formation by ensuring practitioners engage with
problems before encountering AI-generated representations. The pair programming
mindset sustains active cognitive engagement during GenAI-assisted work,
preventing delegation from becoming abdication. These strategies
represent deliberate friction: a structured reintroduction of the cognitive
effort that constitutes architectural expertise.

\subsubsection{The Maturity Model as Diagnostic}

The four-stage maturity model (Table~\ref{tab:maturity-model}) received partial
empirical support through the asynchronously collected SP2 maturity progression
vote (Table~\ref{tab:maturity-progression}). The panel strongly endorsed the natural progression assumption while rejecting
the proposition that short experience can substitute for sustained engagement
(Table~\ref{tab:maturity-progression}). The model's value remains diagnostic rather than prescriptive. Practitioners can
use it to locate their current practice and identify potential growth areas
without implying that progression through all four stages is necessary or
desirable for every organizational context.

The SP2 results reveal a nuanced position that echoes a broader debate in EA
maturity literature. \textcite{kotusev_enterprise_2018} demonstrated that EA
maturity models often prescribe trajectories that bear little resemblance to how
practices actually evolve. The panel's moderate, high-variability agreement on per-category measurement
and context-dependence aligns with this critique: progression is real but not
uniform across practice categories. The high variability across organizational contexts
(particularly in Governance/Compliance and Stakeholder Communication) further
suggests that context-specific factors shape adoption pace more than any
universal progression logic. The specific four-stage structure remains
researcher-proposed, and the asynchronous collection mode introduces
methodological caveats (Section~\ref{subsec:data-limitations}).

%% ================================================================
%% 8.6 Implications for Theory
%% ================================================================
\subsection{Implications for Theory}\label{subsec:disc-implications-theory}

The findings extend three theoretical areas: distributed cognition, practice
theory, and the conceptualization of GenAI in professional knowledge work.

\subsubsection{GenAI as Cognitive Partner}

Distributed cognition theory \parencite{hutchins_cognition_1995} originally
examined how cognitive processes distribute across human team members and
physical artifacts. The Delphi findings suggest that GenAI occupies a position
in the cognitive system distinct from both. The organizational context
boundary, near-unanimously identified as a hard limit on delegation
(Section~\ref{subsubsec:cross-cutting}), delineates where GenAI's
cognitive partnership reaches its current limits. Yet within that boundary,
practitioners described GenAI generating novel representations, responding to
queries contextually, and adapting outputs to conversational framing;
capabilities that exceed those of the charts, instruments, and procedures in
\textcite{hutchins_cognition_1995}'s navigation studies. In our view, GenAI
functions as a new category of cognitive partner: more capable than a static
artifact but less reliable than a human collaborator in domains requiring
tacit organizational knowledge.

These findings point toward a pattern not explicitly addressed in
\textcite{hollan_distributed_2000}'s three types of distribution: distribution
to agents whose reliability varies across task types. The consensus--variability
profiles from Section~\ref{subsec:disc-jagged-frontier} offer preliminary
indications of this pattern. The tight consensus on System/Codebase Analysis suggests relatively stable
perceived reliability, while the wide disagreement on Stakeholder Communication
suggests less predictable performance. These variability differences reflect
panel-level variation in importance ratings rather than direct measures of
GenAI reliability. Yet the
pattern they trace is consistent with the jagged frontier
\parencite{dellacqua_navigating_2023}, where AI capability varies
unpredictably across task boundaries. While the present study's scope does not
permit theoretical generalization, the observation warrants further
investigation. Distributed cognition theory, as originally formulated, assumed
relatively stable capabilities among system components. GenAI may introduce a
partner whose competence shifts across task boundaries, requiring continuous
calibration; a dynamic the original framework did not anticipate.

\subsubsection{Practice Emergence Under Rapid Technological Change}

\textcite{feldman_theorizing_2011}'s practice theory provides a lens for
understanding how recurrent patterns of action emerge and stabilize. The Delphi
findings illustrate a distinctive feature of practice emergence under rapid
technological change: the technology's capabilities shift faster than practices
can stabilize.

The seven practice categories validated in Round~2 showed no challenges or
revisions in Round~3 (Section~\ref{subsubsec:convergence}), suggesting that
the practice taxonomy itself has stabilized. However, the specific
enactments within each category (the tools used, the verification strategies
applied, the interaction patterns adopted) continue to evolve as GenAI
capabilities advance. This creates what we might describe as stable categories
with fluid content: the practice patterns are recognizable but their
instantiation changes continuously.

\textcite{orlikowski_using_2000} argued that technology-in-practice emerges
through situated use rather than following predetermined adoption paths.
The Delphi findings support this claim but add a temporal dimension: when
the technology itself evolves rapidly, practitioners must continuously
re-situate their practices. The frontier learning pattern (the majority of participants discovering
limitations through errors) demonstrates this ongoing re-situation in action.

\subsubsection{The Junior Colleague Mental Model}

The ``junior colleague'' framing, recurrent across the Round~3 panel, deserves
theoretical attention as a novel calibration heuristic. One participant
stated: ``Think of AI as a junior architect and yourself as the senior
architect'' (R3 Brainstorm). Another described treating GenAI interaction like
a buddy system where ``sharing the elements as if it were human betters the
end output'' (R3 Brainstorm).

This mental model is not a formal theoretical contribution but a practitioner
heuristic that performs important cognitive work. It establishes appropriate
expectations (capable but requiring supervision), defines the appropriate
relationship (mentorship rather than delegation or replacement), and provides
a familiar frame of reference for navigating an unfamiliar technology.

In our view, the junior colleague model is theoretically significant because
it reveals how practitioners make sense of a novel cognitive partner by drawing
on social frameworks from human collaboration. This process (making the
unfamiliar familiar through analogical reasoning) represents a sense-making
strategy that existing literature on GenAI adoption has not foregrounded. The
framing also reveals its own limitations. Unlike a junior colleague, GenAI
does not learn from supervision in ways that transfer across sessions. Its
failure modes also differ fundamentally from those of a human novice.

%% ================================================================
%% 8.7 Limitations
%% ================================================================
\subsection{Limitations}\label{subsec:disc-limitations}

Several limitations qualify the findings and their interpretation.

\textbf{Sample size and composition.} The Delphi panel comprised eight to ten
participants across three rounds, sufficient for a Delphi study targeting information richness but limiting
generalizability. The panel was
predominantly European, self-selected through the researcher's professional
network, and composed of practitioners already using GenAI. This composition
produces a portrait of early adopter practices that may not represent the
broader EA practitioner population.

\textbf{Scope boundary tension.} The workbook was framed as ``Software
Architecture / Solution Design work'' rather than ``Enterprise Architecture'' to
capture practitioners who do RA work under various professional labels (see
Section~\ref{sec:methodology}). This terminological broadening may have
admitted practices that sit at the boundary between RA-specific and general
architectural work. The workbook's RA lifecycle mapping (Part~C) and R2's
RA-contextualized validation mitigate this risk, but some practices (e.g.,
System/Codebase Analysis) may represent general architectural work that is not
exclusive to RA development. This is an inherent trade-off: ecological validity
(capturing how practitioners actually describe their work) versus strict
construct alignment with the research question's ``EA practitioners'' scope.
Future research could isolate RA-specific practices from general architectural
practices through comparative studies across different architectural artifact
types.

\textbf{Temporal snapshot.} Data collection occurred between January and March
2026, during a period of rapid GenAI capability evolution. Findings represent
a snapshot of practice at a specific moment. GenAI capabilities available
during the study period (including large context windows, multi-modal
reasoning, and agent-based workflows) may shift substantially within months,
potentially rendering specific findings (particularly tool-dependent practices
like Deep Workflow Integration) outdated.

\textbf{Partially validated maturity model.} The four-stage maturity model
(Table~\ref{tab:maturity-model}) was derived from observed variation in the
Delphi data. The asynchronously collected SP2 vote
(Table~\ref{tab:maturity-progression}) provides empirical support for the underlying progression assumption
($M = 4.71$, variability 22\%), partially grounding a maturity model within a
consensus-based study. However, a residual design-science tension remains: the specific
four-stage structure (from Experimental to Organizational Embedding) was not
individually validated by the panel, the vote was collected asynchronously
without preceding group discussion, and the sample comprised seven participants.
Future research should validate the specific stage boundaries and
characteristics.

\textbf{Delphi process constraints.} Three rounds may limit the depth of
convergence achievable. Round~1's asynchronous format allowed thoughtful
individual responses but prevented the real-time probing possible in
interviews. Rounds~2 and~3 compressed discussion into 120-minute sessions,
with Round~3 Step~15 not executed and Step~16 collected asynchronously. The
absence of Step~15 (Future Perspectives) means the temporal dimension and
future-oriented perspectives of the conceptual framework remain empirically
underexplored.

\textbf{Researcher positionality.} The researcher is an EA practitioner who
uses GenAI daily for architectural work. This dual role means the researcher
is simultaneously a research subject and the person designing the inquiry.
Insider status touched every phase of the research: the discovery workbook
drew on practitioner knowledge of EA workflows, coding decisions reflected
familiarity with architectural reasoning patterns, the seven practice
categories emerged through researcher synthesis of participant data, and the
guidance artifact incorporated the researcher's understanding of what
practitioners would find actionable. Each of these steps involved interpretive
judgments that insider knowledge could shape, creating a structural bias toward
validating the significance of GenAI integration.

Assumption management relied primarily on the Delphi method's structural
safeguards rather than a rigorous personal reflexivity protocol. We acknowledge
that the researcher's actual reflexive practice was less systematic than
planned; reflexive notes were captured at key decision points rather than
maintained as a continuous journal. Three structural features partially
compensated for this gap. Round~1's asynchronous workbook format ensured
participants articulated practices in their own terms before researcher coding
began. Round~2's live voting and discussion allowed the panel to validate or
reject practice categories collectively, not through researcher judgment alone.
Round~3's structured brainstorming generated participant responses that were
aggregated before researcher synthesis. In our view, these deliberative
safeguards provided meaningful protection against insider bias, but we
acknowledge that a more disciplined personal reflexivity practice would have
strengthened confirmability.

Two reflexive insights merit disclosure. The researcher has experimented with
large language models since 2017 and assumed that most EA practitioners would
gradually incorporate GenAI over time. The genuine surprise was the speed of
conversion: several participants described themselves as recently skeptical,
believing EA work was too context-heavy for AI assistance, yet had rapidly
become active daily users within months. This challenged the researcher's
gradualist assumption and highlighted the threshold-crossing pattern visible in
the findings. Regarding the ``Ghost
Architect'' provocation, this stimulus emerged primarily from the data rather
than insider projection. Rounds~1 and~2 surfaced unanticipated cross-role
dynamics, including developers producing architectural artifacts bottom-up
(Section~\ref{subsubsec:cross-role}). The provocation was designed to probe
these participant-identified tensions in Round~3. We acknowledge that insider
status likely influenced the salience assigned to role boundary shifts; an
active practitioner may notice such dynamics more readily. However, the
provocation's content was anchored in participant-generated themes rather than
researcher preconceptions.

\textbf{GenAI tool dependency.} Findings are shaped by the GenAI tools
available in 2025--2026, including large language models from major providers,
IDE-integrated coding assistants, and early agent-based workflows. The
practices, heuristics, and governance considerations identified may not
transfer directly to future GenAI paradigms with different capability profiles.

%% ================================================================
%% 8.8 Future Research
%% ================================================================
\subsection{Future Research}\label{subsec:future-research}

Five directions merit investigation. First, longitudinal research should track
how the seven practice categories evolve as GenAI capabilities advance, since
the temporal snapshot limitation means that practices documented here may
fragment, merge, or transform as tool capabilities shift.

Second, cross-organizational validation of the practice taxonomy would establish
whether the categories identified in this early-adopter panel hold across
different organizational contexts, governance maturity levels, and industry
sectors. Replication studies with larger panels would test whether the consensus
patterns observed here generalize beyond this exploratory sample.

Third, the maturity model's progression assumption received strong panel support
($M = 4.71$), but three areas warrant further validation: the specific four-stage
structure (from Experimental to Organizational Embedding) was not individually
tested, the practice-category-specific nature of maturity suggested by the panel
implies different progression trajectories per practice may exist, and the
model's applicability beyond the early-adopter population studied here remains
unknown.

Fourth, sector-specific RA contexts (healthcare, financial services, government)
introduce regulatory and compliance dimensions that may reshape which
practice categories dominate and how verification heuristics are calibrated.

Fifth, the anti-atrophy strategies proposed in
Section~\ref{subsec:disc-implications-practice} require longitudinal validation.
While this study identifies the competency decay paradox and proposes deliberate
friction interventions, their long-term effectiveness in preserving
architectural expertise remains an open empirical question with significant
implications for EA education and professional development.

%% ================================================================
%% 8.9 Closing Statement
%% ================================================================
\subsection{Closing Statement}\label{subsec:closing-statement}

The research began with a practice vacuum: EA practitioners were integrating
GenAI into RA development without empirical evidence, evaluation frameworks,
shared vocabulary, or validated guidance.
\textcite{kotusev_enterprise_2018} called for reconceptualizing EA practice to
reflect how architects actually work; GenAI has made that call urgent. Through
three Delphi rounds, this study maps what practitioners are doing, how their
work is changing, and what guidance helps. The practices documented here will
evolve as GenAI capabilities advance. What persists is the underlying question:
which parts of architectural work to delegate, and which to keep.
